% GENERATED FILE. Edit canonical lessons, then run npm run generate:books.
% canonical-source-hash: fnv1a64:118aafe66fe48dc9
% canonical-lessons: FA-C296-khoshhal-budam, FA-C296-ghaza, FA-W296-gheyn, FA-C296-omidvaram-ke, FA-R296-feelings-and-hopes

\chapter{Glad, Tired, Hopeful}
\label{ch:fa-glad-tired-hopeful}
\hlchaptermodality{296}

\begin{chapteropening}
\textbf{By the end of this chapter:} \emph{I can say how I felt and what I became, describe a meal, and say what I hope for and would like to do.}

Complete FA-R296-feelings-and-hopes: how a day made you feel, and one thing you hope for.
\end{chapteropening}

\section[\texorpdfstring{\fa{خوشحال} \fa{بودم} (khoshhâl budam)}{خوشحال بودم (khoshhâl budam)}]{\fa{خوشحال} \fa{بودم} (khoshhâl budam) — I was happy, and I became happy}

\label{lesson:FA-C296-khoshhal-budam}



\begin{quote}
Say \textquotedblleft{}happy\textquotedblright{}, \textquotedblleft{}tired\textquotedblright{}, \textquotedblleft{}upset\textquotedblright{}. (\emph{Khoshhâl}, \emph{khaste}, \emph{nârâhat}.) Then \textquotedblleft{}to become\textquotedblright{}. (\emph{Shodan.})

You can say what happened on a trip. Describing it means saying how it felt, and Persian gives a feeling two pasts as well.
\end{quote}

\begin{grammarlens}[title={Grammar Lens: the feeling as scene}]
The adjective, then \emph{budan} in the past:

\begin{quote}
\textbf{\fa{خوشحال} \fa{بودم}.} — \emph{khoshhâl budam} — I was happy.
\end{quote}

\emph{Khaste budam}, I was tired; \emph{nârâhat budim}, we were upset; \emph{shâd bud}, she was cheerful. The adjective stays the same for every person, and the ending rides on \emph{bud-}. This is the feeling as background, there for a while.
\end{grammarlens}

\begin{grammarlens}[title={Grammar Lens: the feeling as event}]
Swap \emph{budan} for \emph{shodan}, \textquotedblleft{}to become\textquotedblright{}, and the feeling arrives at a moment:

\begin{quote}
\textbf{\fa{خسته} \fa{شدم}.} — \emph{khaste shodam} — I got tired.
\end{quote}

\begin{quote}
\textbf{\fa{وقتی} \fa{قطار} \fa{را} \fa{از} \fa{دست} \fa{دادم،} \fa{ناراحت} \fa{شدم}.} — \emph{vaqti qatâr râ az dast dâdam, nârâhat shodam} — When I missed the train, I got upset.
\end{quote}

The cause can come with \emph{az}, \textquotedblleft{}from\textquotedblright{}, and the infinitive works as a noun:

\begin{quote}
\textbf{\fa{از} \fa{دیدن} \fa{شما} \fa{خوشحال} \fa{شدم}.} — \emph{az didan-e shomâ khoshhâl shodam} — I was glad to see you.
\end{quote}

Word for word, \textquotedblleft{}from the seeing of you I became happy\textquotedblright{}. It is a thing people say at the end of a first meeting.

So \emph{khoshhâl budam} is the mood of the whole day, and \emph{khoshhâl shodam} is the moment the news arrived.
\end{grammarlens}

\begin{cousinweb}
\emph{Shodan} is usually traced to Middle Persian \emph{šudan}, \textquotedblleft{}to go\textquotedblright{}. Becoming is going into a state: \emph{khaste shodam} is \textquotedblleft{}I went tired\textquotedblright{}. \emph{Khoshhâl} is \emph{khosh}, \textquotedblleft{}good\textquotedblright{}, with \emph{hâl}, \textquotedblleft{}state\textquotedblright{}, the same \emph{hâl} as in \emph{hâl-e shomâ chetor ast?}
\end{cousinweb}

\subsection*{Your turn}
Say these aloud:

\begin{itemize}
\raggedright
  \item \textquotedblleft{}I was happy\textquotedblright{}, then \textquotedblleft{}I was tired\textquotedblright{} — \emph{khoshhâl budam}, \emph{khaste budam}
  \item \textquotedblleft{}I got tired\textquotedblright{} — \emph{khaste shodam}
  \item \emph{vaqti qatâr râ az dast dâdam, nârâhat shodam}
  \item at the end of a first meeting — \emph{az didan-e shomâ khoshhâl shodam}
\end{itemize}

\begin{tcolorbox}[breakable,colback=teal!4,colframe=teal!35!black,title={Before you move on}]
Which is the mood of the day, \emph{khoshhâl budam} or \emph{khoshhâl shodam}? (\emph{\textbf{Budam}}: scene.) Which is the moment? (\emph{\textbf{Shodam}}: event.) \textquotedblleft{}I was glad to see you\textquotedblright{}? (\emph{\textbf{Az didan-e shomâ khoshhâl shodam.}})
\end{tcolorbox}
\section[\texorpdfstring{\fa{غذا} (ghazâ)}{غذا (ghazâ)}]{\fa{غذا} (ghazâ) — food, and how it was}

\label{lesson:FA-C296-ghaza}



\begin{quote}
Say \textquotedblleft{}I was happy\textquotedblright{}, then \textquotedblleft{}I got tired\textquotedblright{}. (\emph{Khoshhâl budam}, \emph{khaste shodam}.) Then \textquotedblleft{}delicious\textquotedblright{}. (\emph{Khoshmaze.})

Talk about a trip, and sooner or later the food comes up.
\end{quote}

\subsection*{You'll want to know: \fa{غذا}}
\textbf{\fa{غذا}} — \emph{ghazâ} — \textquotedblleft{}food\textquotedblright{}, and \textquotedblleft{}a meal\textquotedblright{}.

\begin{quote}
\textbf{\fa{غذای} \fa{این} \fa{رستوران} \fa{خیلی} \fa{خوشمزه} \fa{بود}.} — \emph{ghazâ-ye in resturân kheyli khoshmaze bud} — The food at this restaurant was very tasty.
\end{quote}

After its long \emph{â} the ezafe is \emph{-ye}, and here it is written, as \textbf{\fa{ی}}: \textbf{\fa{غذای}}. \emph{Ghazâ khordan} is \textquotedblleft{}to eat, to have a meal\textquotedblright{}: \emph{ghazâ khordim}, we ate.

Say \emph{gha-zâ}, with the \emph{gh} far back, like a soft gargle. Its first letter is new, and the next lesson puts it in your hand. The second is \textbf{\fa{ذ}}, which you wrote in chapter two hundred and ninety.

\begin{cousinweb}
\emph{ghazâ} is borrowed from Arabic \emph{ghidhāʾ}, \textquotedblleft{}nourishment, food\textquotedblright{} (EtymDB 2.1). Arabic says its middle letter as the \emph{th} of English \emph{this}. Persian kept the letter \textbf{\fa{ذ}} in the spelling and says it \emph{z}, which is why \emph{ghazâ} sounds nothing like the Arabic it is spelled from.
\end{cousinweb}

\subsection*{Your turn}
Say these aloud:

\begin{itemize}
\raggedright
  \item \emph{ghazâ}, with the \emph{gh} far back
  \item \emph{ghazâ-ye in resturân kheyli khoshmaze bud}
  \item \textquotedblleft{}we ate\textquotedblright{}, with the noun — \emph{ghazâ khordim}
  \item an opinion of a meal — \emph{be nazar-e man ghazâ khoshmaze bud}
\end{itemize}

\begin{tcolorbox}[breakable,colback=teal!4,colframe=teal!35!black,title={Before you move on}]
\textquotedblleft{}Food\textquotedblright{}? (\emph{\textbf{Ghazâ.}}) How is the ezafe written after it? (\emph{\textbf{-ye}}, as \textbf{\fa{ی}}: \textbf{\fa{غذای}}.) Which letter of \textbf{\fa{غذا}} sounds \emph{z}? (\textbf{\fa{ذ}}, \emph{zâl}.)
\end{tcolorbox}
\section[\texorpdfstring{\fa{غ} (gheyn)}{غ (gheyn)}]{\fa{غ} (gheyn) — the letter that opens \fa{غذا}}

\label{lesson:FA-W296-gheyn}



\begin{quote}
Before the new one: write \fa{ع} and \fa{ق} once each, and name them.

Now say \textbf{\fa{غذا}} (\emph{ghazâ}), \textquotedblleft{}food\textquotedblright{}. You met it a moment ago. This lesson takes its first letter out and puts it in your hand.
\end{quote}

\subsection*{Script you'll notice: \fa{غ}}
\textbf{\fa{غ}} — \emph{gheyn}, the sound \emph{gh}, made far back, like a soft gargle.

It opens \textbf{\fa{غذا}}, in its opening shape, \textbf{\fa{غـ}}, joined to \textbf{\fa{ذ}}. It is \textbf{\fa{ع}} with \textbf{one dot above}, and it changes shape the same way: an open hook at the start, a small closed loop in the middle, a tail at the end.

In Tehran speech \textbf{\fa{غ}} and \textbf{\fa{ق}} sound alike: \emph{ghazâ} and \emph{qatâr} open on the same sound. Spelling keeps them apart, and so does this book's romanization: \textbf{\fa{ق}} is written \emph{q}, as in \emph{qatâr}, and \textbf{\fa{غ}} is written \emph{gh}, as in \emph{ghazâ}. Each word has to be learned with its own letter.

\subsection*{Writing: \fa{غ}}
Put your pen on \fa{غ} and follow its line. Copy the shape you can see — slowly, and larger than it is printed.

\begin{quote}
This book does not yet tell you \textbf{where to start this letter or which way to travel}. It is not written down here until it can be written down with a source. Copying what is in front of you needs no such source, and it is how the shape gets into your hand in the meantime.
\end{quote}

\subsection*{Your turn}
\begin{itemize}
\raggedright
  \item \emph{Find:} \fa{غ} at the start of \textbf{\fa{غذا}}
  \item \emph{Trace it:} \fa{غ} standing alone three times, saying \emph{gh} as you finish each one
  \item \emph{Write it:} \fa{ع}, then \fa{غ}, from memory: one dot apart
  \item \emph{Write it:} \textbf{\fa{غذا}}, then \textbf{\fa{قطار}}: one sound, two letters
  \item \emph{Read it:} \textbf{\fa{غذا}} aloud once more
\end{itemize}

\begin{tcolorbox}[breakable,colback=teal!4,colframe=teal!35!black,title={Before you move on}]
Which letter is this — \fa{غ}? (\emph{\textbf{gheyn}}, \emph{gh}.) Which word did it come from? (\textbf{\fa{غذا}}, \emph{ghazâ}, \textquotedblleft{}food\textquotedblright{}.) Which other letter sounds the same in Tehran? (\textbf{\fa{ق}}, written \emph{q}.)
\end{tcolorbox}
\section[\texorpdfstring{\fa{امیدوارم} \fa{که} (omidvâram ke)}{امیدوارم که (omidvâram ke)}]{\fa{امیدوارم} \fa{که} (omidvâram ke) — I hope that, I'd like to, I dream of}

\label{lesson:FA-C296-omidvaram-ke}



\begin{quote}
Say \textquotedblleft{}to hope\textquotedblright{}, then \textquotedblleft{}I want to go\textquotedblright{}. (\emph{Omidvâr budan}, \emph{mikhâham beravam}.) Then \textquotedblleft{}I'm well\textquotedblright{}, in one word. (\emph{Khubam.})

Describing an experience runs forward as well as back: what you hope for next, and what you would like to become.
\end{quote}

\begin{grammarlens}[title={Grammar Lens: omidvâram ke, and the subjunctive}]
\emph{Omidvâr} takes the \emph{-am} of \emph{khubam}: \textbf{\fa{امیدوارم}}, \emph{omidvâram}, \textquotedblleft{}I am hopeful\textquotedblright{}, \textquotedblleft{}I hope\textquotedblright{}. What you hope for follows \emph{ke} in the subjunctive, because it has not happened:

\begin{quote}
\textbf{\fa{امیدوارم} \fa{که} \fa{سال} \fa{دیگر} \fa{به} \fa{تهران} \fa{بروم}.} — \emph{omidvâram ke sâl-e digar be Tehrân beravam} — I hope to go to Tehran next year.
\end{quote}

Chapter one hundred and thirty-eight showed \emph{omidvâram hâlat khub bâshad}, \textquotedblleft{}I hope you are well\textquotedblright{}. The \emph{bâshad} there is this subjunctive, and the \emph{ke} has been left out, as speech often does.
\end{grammarlens}

\begin{grammarlens}[title={Grammar Lens: what you'd like, and what you dream of}]
\emph{Dust dâram}, \textquotedblleft{}I like\textquotedblright{}, with a subjunctive after it is \textquotedblleft{}I'd like to\textquotedblright{}:

\begin{quote}
\textbf{\fa{دوست} \fa{دارم} \fa{مهندس} \fa{بشوم}.} — \emph{dust dâram mohandes beshavam} — I'd like to become an engineer.
\end{quote}

\emph{Beshavam} is \emph{shodan} in the subjunctive, on its stem \emph{shav-}; in speech it is shortened to \emph{besham}. \emph{Ârezu}, \textquotedblleft{}a wish, a dream\textquotedblright{}, works the same way with \emph{dâram}:

\begin{quote}
\textbf{\fa{آرزو} \fa{دارم} \fa{که} \fa{یک} \fa{روز} \fa{خلبان} \fa{بشوم}.} — \emph{ârezu dâram ke yek ruz khalabân beshavam} — My dream is to become a pilot one day.
\end{quote}

So three strengths: \emph{mikhâham}, I want; \emph{dust dâram}, I'd like; \emph{ârezu dâram}, I dream of. All three take the subjunctive.
\end{grammarlens}

\subsection*{Your turn}
Say these aloud:

\begin{itemize}
\raggedright
  \item \emph{omidvâram ke sâl-e digar be Tehrân beravam}
  \item \textquotedblleft{}I'd like to become an engineer\textquotedblright{} — \emph{dust dâram mohandes beshavam}
  \item a dream of your own, with \emph{ârezu dâram ke} and \emph{beshavam}
  \item a meal, then a hope — \emph{ghazâ khoshmaze bud. omidvâram ke dobâre beravim}
\end{itemize}

\begin{tcolorbox}[breakable,colback=teal!4,colframe=teal!35!black,title={Before you move on}]
What is \emph{omidvâram}, taken apart? (\emph{Omidvâr} with \emph{\textbf{-am}}, \textquotedblleft{}I am\textquotedblright{}.) Why is the verb after it in the subjunctive? (The hoped-for thing \textbf{has not happened}.) Which letter opens \textbf{\fa{غذا}}? (\textbf{\fa{غ}}, \emph{gheyn}.)
\end{tcolorbox}
\section[(review)]{Review — how it felt, and what you hope for}

\label{lesson:FA-R296-feelings-and-hopes}



\begin{quote}
Say \textquotedblleft{}I was tired\textquotedblright{}, then \textquotedblleft{}I got tired\textquotedblright{}. (\emph{Khaste budam}, \emph{khaste shodam}.)
\end{quote}

\begin{grammarlens}[title={What you've built}]
Three lessons, and an experience now has its colour. \emph{Budam} gives the feeling of a whole stretch; \emph{shodam} gives the moment it arrived. \emph{Ghazâ} brings in the meal, and \emph{omidvâram ke}, \emph{dust dâram} and \emph{ârezu dâram} carry the hope forward, each with the subjunctive after it.

Back to the bad morning at the station, told as an experience:

\begin{quote}
\emph{qatâr râ az dast dâdam va nârâhat shodam. ba'd mas'ul-e istgâh bilit-e fardâ râ be man dâd, barâ-ye hamin khoshhâl shodam. ghazâ-ye istgâh khoshmaze bud. omidvâram ke fardâ dir naresam.}
\end{quote}

The last verb, \emph{naresam}, is the subjunctive with \emph{na-}: \textquotedblleft{}that I don't arrive\textquotedblright{}.
\end{grammarlens}

\subsection*{Your turn}
Say these aloud:

\begin{itemize}
\raggedright
  \item the station morning, told as an experience, without stopping
  \item a whole day as mood — \emph{ruz khub bud}, the day was good, then \emph{khoshhâl budam}
  \item something you'd like, and something you dream of — \emph{dust dâram...}, \emph{ârezu dâram ke...}
  \item the hope at the end — \emph{omidvâram ke fardâ dir naresam}
\end{itemize}

\begin{tcolorbox}[breakable,colback=teal!4,colframe=teal!35!black,title={Before you move on}]
Which tells the mood of a day, \emph{budam} or \emph{shodam}? (\emph{\textbf{Budam.}}) \textquotedblleft{}I hope I won't be late tomorrow\textquotedblright{}? (\emph{\textbf{Omidvâram ke fardâ dir naresam.}}) Which two letters of this run are new to your hand? (\textbf{\fa{غ}} in \textbf{\fa{غذا}}, \textbf{\fa{ئ}} in \textbf{\fa{مسئول}}.)
\end{tcolorbox}
